% !TEX program = lualatex
%!TEX TS-program = lualatex
\documentclass[12pt,a4paper]{article}
\usepackage{fontspec}
\usepackage{polyglossia}
\setdefaultlanguage{polish}
\usepackage{geometry}
\geometry{margin=2.1cm}
\usepackage[table]{xcolor}
\usepackage{longtable}
\usepackage{array}
\usepackage{booktabs}
\usepackage{enumitem}
\usepackage{ragged2e}
\definecolor{tableHeader}{HTML}{E8EEF7}
\definecolor{tableStripe}{HTML}{F7F9FC}
\newcolumntype{L}[1]{>{\RaggedRight\arraybackslash}p{#1}}
\newcolumntype{C}[1]{>{\Centering\arraybackslash}p{#1}}
\renewcommand{\arraystretch}{1.22}
\setlength{\tabcolsep}{4pt}
\setlist[itemize]{leftmargin=*, itemsep=2pt, topsep=2pt}
\setlength{\parindent}{0pt}
\setlength{\parskip}{6pt}
\emergencystretch=3em
\sloppy
\begin{document}

\section*{Raport ze sprintu nr 6 - Calisia Banking App}
Data raportu: 31.05.2026

Okres sprintu: 26.05 - 31.05.2026

Product Owner / Scrum Master: Dawid Olszewski

\subsection*{Lista osób}
\begin{itemize}
\item Kubryn Jeremiasz
\item Łaszuk Jakub
\item Makarevich Ksenia
\item Nuckowski Michał
\item Ruczyńska Nikola
\item Wołczko Jakub
\item Olszewski Dawid
\end{itemize}

\subsection*{Product Owner i Scrum Master sprintu}
Dawid Olszewski

\subsection*{Lista historyjek na sprint}
\begingroup
\footnotesize
\rowcolors{2}{tableStripe}{white}
\begin{longtable}{@{}C{0.08\textwidth}L{0.62\textwidth}C{0.10\textwidth}L{0.12\textwidth}@{}}
\toprule
\rowcolor{tableHeader}
\textbf{Kod} & \textbf{Historyjka} & \textbf{SP} & \textbf{Status} \\
\midrule
\endfirsthead
\toprule
\rowcolor{tableHeader}
\textbf{Kod} & \textbf{Historyjka} & \textbf{SP} & \textbf{Status} \\
\midrule
\endhead
S6-01 & Przelew własny między rachunkami klienta & 13 & Done \\
S6-02 & Przelew zewnętrzny & 13 & Done \\
S6-03 & Panel przelewów i szybka aktualizacja dashboardu & 8 & Done \\
S6-04 & Walidacje i scenariusze negatywne przelewów & 8 & Done \\
\bottomrule
\end{longtable}
\endgroup

\subsection*{Podsumowanie sprintu}
Cel sprintu: Przelewy własne i zewnętrzne. Pozwolenie klientowi na wykonywanie przelewów oraz odzwierciedlanie ich w saldach i historii.

Zakres sprintu obejmował 4 historyjek o łącznej estymacji 42 SP. Wszystkie historyjki zostały dowiezione ze statusem Done, dlatego osiągnięta prędkość sprintu wyniosła 42 SP.

Dostarczono obsługę przelewów własnych i zewnętrznych, w tym walidacje danych, obciążenie rachunku źródłowego, zapis transferu oraz aktualizację dashboardu po operacji.

Sprint domknął najważniejszą funkcjonalność operacyjną bankowości. Klient może wykonać przelew, a system utrzymuje spójność sald, historii i widoku dashboardu.

\subsubsection*{Najważniejsze rezultaty}
\begin{itemize}
\item Dodano przelew własny między rachunkami klienta.
\item Dodano przelew zewnętrzny na rachunek odbiorcy.
\item Połączono panel przelewów z aktualizacją dashboardu.
\item Przetestowano scenariusze negatywne i walidacje przelewów.
\end{itemize}

\subsection*{Retrospekcja sprintu}
Retrospekcja potwierdziła, że funkcje finansowe wymagają szczególnie precyzyjnych reguł domenowych i dobrze opisanych scenariuszy testowych.

\subsubsection*{Co się udało}
\begin{itemize}
\item Udało się zamodelować przelewy własne i zewnętrzne.
\item Frontend zapewnił panel przelewów z szybkim odświeżeniem dashboardu.
\item Testy pokryły brak środków, błędne dane i nieautoryzowany dostęp.
\item Kolekcja Bruno umożliwiła szybkie sprawdzanie przepływu API.
\end{itemize}
\subsubsection*{Poziom energii zespołu}
Średni poziom energii zespołu: 4,0/5.

Energia była dobra, choć sprint wymagał skupienia ze względu na złożone reguły finansowe i wiele scenariuszy negatywnych.

\subsubsection*{Czego udało się nauczyć}
\begin{itemize}
\item Zespół nauczył się dokładniej rozdzielać przelewy własne i zewnętrzne.
\item Scenariusze negatywne są kluczowe przy funkcjach finansowych.
\item Aktualizacja dashboardu po akcji użytkownika powinna być projektowana razem z API.
\end{itemize}
\subsubsection*{Wnioski i działania na kolejny sprint}
\begin{itemize}
\item Przygotować checklistę regresji dla pełnego przepływu finansowego.
\item W stabilizacji położyć nacisk na spójność komunikatów błędów.
\item Uzupełnić scenariusz demo o przelew i odświeżenie dashboardu.
\end{itemize}

\subsection*{Następny sprint}
Product Owner / Scrum Master w następnym sprincie: Nikola Ruczyńska

Główny cel następnego sprintu: Stabilizacja, testy i oddanie produktu. Dopracowanie całego przepływu klienta, usunięcie niespójności i przygotowanie produktu do prezentacji.

\end{document}
